\documentclass[11pt]{article}

\usepackage[margin=1in]{geometry}
\usepackage{amsmath}
\usepackage{amssymb}
\usepackage{booktabs}
\usepackage{longtable}
\usepackage{array}
\usepackage{xcolor}
\usepackage[hidelinks]{hyperref}

\title{Functional Dependencies and Minimal Basis}
\author{}
\date{}

\newcommand{\fd}{\rightarrow}
\newcommand{\code}[1]{\texttt{#1}}

\begin{document}

\maketitle

\section{Purpose}

This document records the functional dependency identification and minimal-basis
work for the Sprint Two goal \emph{Calculate Minimal Bases and Project Relational
Dependencies}. The projection and dependency-preservation portions of the goal
will be documented separately. This analysis is based on the current Sprint Two
ERD.


\section{Updated Relational Schema Used for the Analysis}

The updated ERD maps to the following relations:

\begin{align*}
\text{Region}(&region\_id, region\_name, country)\\
\text{GPUType}(&gpu\_type\_name)\\
\text{SpecificGPUType}(&gpu\_type\_name, memory\_gb)\\
\text{AIModel}(&model\_id, model\_name)\\
\text{GPUServer}(&machine\_id, cpu\_capacity, memory\_capacity, gpu\_capacity,\\
&region\_id, gpu\_type\_name)\\
\text{InferenceRequest}(&request\_id, submitted\_at, input\_size, status,\\
&response\_time\_ms, region\_id, model\_id, machine\_id)\\
\text{MachineMetric}(&machine\_id, metric\_id, start\_time, end\_time,\\
&worker\_load, gpu\_utilization, cpu\_utilization)\\
\text{Hosts}(&machine\_id, model\_id)\\
\text{CompatibleWith}(&model\_id, gpu\_type\_name)
\end{align*}

The primary key of \code{MachineMetric} is now $(machine\_id, metric\_id)$ because
\code{MachineMetric} is a weak entity owned by \code{GPUServer}. The key of
\code{SpecificGPUType} is \code{gpu\_type\_name}, it is also a foreign key to the
parent \code{GPUType} relation.

\subsection*{Region role notation}

The ERD has two different relationships to \code{Region}: a request originates in
a region, while a server is located in a region. Both local SQL columns are named
\code{region\_id}, but they represent different roles. To avoid accidentally
identifying them as the same attribute during the FD calculations, this document
uses the temporary names

\[
request\_region\_id \qquad\text{and}\qquad server\_region\_id.
\]

When projected back onto the individual relations, both are written as the local name \code{region\_id}.

\section{Functional Dependencies from the ERD}

The identifier of a strong entity
functionally determines its non-identifier attributes. A many-to-one relationship
also means the identifier on the many side determines the identifier on the one
side. For the weak \code{MachineMetric} entity, the full identifier is the union of
the owner identifier and the weak entity's local identifier. Before splitting right-hand sides, the non-trivial project dependencies are:

\begin{align*}
region\_id &\fd \{region\_name, country\}\\
model\_id &\fd model\_name\\
gpu\_type\_name &\fd memory\_gb \quad \text{(within SpecificGPUType)}\\
machine\_id &\fd \{cpu\_capacity, memory\_capacity, gpu\_capacity,\\
&\qquad server\_region\_id, gpu\_type\_name\}\\
request\_id &\fd \{submitted\_at, input\_size, status, response\_time\_ms,\\
&\qquad request\_region\_id, model\_id, machine\_id\}\\
\{machine\_id, metric\_id\} &\fd \{start\_time, end\_time, worker\_load,\\
&\qquad gpu\_utilization, cpu\_utilization\}.
\end{align*}

\code{Hosts} and \code{CompatibleWith} are many-to-many relations with no
non-key relationship attributes. Therefore neither relation contributes a
non-trivial FD between its two participating identifiers. Their composite key
constraints still remain part of the relational schema.

\section{Minimal Basis}
\begin{enumerate}
    \item Split every FD so that each right-hand side contains one attribute.
    \item Try removing each FD, if the remaining FDs still imply it, remove it.
    \item Try removing attributes from the left-hand side of each FD; if the remaining left-hand side still implies the right-hand side, we will remove the extraneous attribute.
\end{enumerate}

\subsection{Step 1: Split the right-hand sides}

After using the splitting rule, the candidate basis is:

\begin{align*}
region\_id &\fd region\_name\\
region\_id &\fd country\\
model\_id &\fd model\_name\\
gpu\_type\_name &\fd memory\_gb\\
machine\_id &\fd cpu\_capacity\\
machine\_id &\fd memory\_capacity\\
machine\_id &\fd gpu\_capacity\\
machine\_id &\fd server\_region\_id\\
machine\_id &\fd gpu\_type\_name\\
request\_id &\fd submitted\_at\\
request\_id &\fd input\_size\\
request\_id &\fd status\\
request\_id &\fd response\_time\_ms\\
request\_id &\fd request\_region\_id\\
request\_id &\fd model\_id\\
request\_id &\fd machine\_id\\
\{machine\_id, metric\_id\} &\fd start\_time\\
\{machine\_id, metric\_id\} &\fd end\_time\\
\{machine\_id, metric\_id\} &\fd worker\_load\\
\{machine\_id, metric\_id\} &\fd gpu\_utilization\\
\{machine\_id, metric\_id\} &\fd cpu\_utilization.
\end{align*}

\subsection{Step 2: Remove redundant FDs}

Each FD was tested by temporarily removing it and calculating the closure of its left-hand side under the remaining dependencies. None of the removed right-hand-side attributes could still be derived, so no redundant FD was found.

For example, if
\[
machine\_id \fd gpu\_type\_name
\]
is removed, $\{machine\_id\}^{+}$ under the remaining dependencies does not
contain \code{gpu\_type\_name}. Therefore that FD must remain.

\subsection{Step 3: Remove extraneous left-hand-side attributes}

All determinants are single attributes except the weak entity dependencies for
\code{MachineMetric}. For those dependencies:

\begin{itemize}
    \item $\{metric\_id\}^{+}$ does not determine any metric measurement
          attributes, so \code{machine\_id} cannot be removed.
    \item $\{machine\_id\}^{+}$ determines server information, but it does not
          determine a particular metric record's measurement attributes, so
          \code{metric\_id} cannot be removed.
\end{itemize}

Therefore $\{machine\_id, metric\_id\}$ is minimal on the left-hand side.

\subsection{Conclusion}

After applying all three minimal-basis steps from lecture, the candidate basis from Section 4.1 satisfies the required conditions:

\begin{itemize}
    \item every functional dependency has a single attribute on the right-hand side;
    \item removing any one of the listed dependencies prevents its right-hand-side
          attribute from being derived from the same determinant; and
    \item no attribute can be removed from the composite determinant
          $\{machine\_id, metric\_id\}$ without losing the corresponding
          \code{MachineMetric} dependency.
\end{itemize}

Thus no dependency is redundant and no determinant contains an extraneous
attribute. The split dependency set listed in Section 4.1 is the project's
minimal basis, denoted $F_{min}$. Projection of $F_{min}$ onto the final relations will be handled in a separate document.

\end{document}
